%%
%% ICML 2025 LaTeX Template
%% Template for International Conference on Machine Learning submissions
%%

\documentclass{article}

% ICML Packages
\usepackage{icml2025}  % ICML style file (needs to be downloaded separately)
\usepackage[T1]{fontenc}
\usepackage{type1cm}
\usepackage{hyperref}
\usepackage{url}
\usepackage{booktabs}
\usepackage{amsfonts}
\usepackage{amsmath}
\usepackage{amssymb}
\usepackage{nicefrac}
\usepackage{microtype}
\usepackage{graphicx}
\usepackage{algorithm}
\usepackage{algorithmic}
\usepackage{xcolor}

% For hyperlinks
\hypersetup{
    colorlinks=true,
    linkcolor=blue,
    citecolor=blue,
    urlcolor=blue
}

% Math operators
\DeclareMathOperator*{\argmin}{arg\,min}
\DeclareMathOperator*{\argmax}{arg\,max}

% Title and authors
\icmltitlerunning{Short Title for Running Header}

\begin{document}

\twocolumn[
\icmltitle{Your Paper Title Here\\
           With Possible Subtitle}

\icmlsetsymbol{equal}{*}

\begin{icmlauthorlist}
\icmlauthor{First Author\icmlequal{} }{aff1,aff2}
\icmlauthor{Second Author}{aff1}
\icmlauthor{Third Author}{aff2}
\end{icmlauthorlist}

\icmlaffiliation{aff1}{Department of Computer Science, University A}
\icmlaffiliation{aff2}{Department of Machine Learning, University B}

\icmlcorrespondingauthor{First Author}{first.author@email.com}

\icmlkeywords{Machine Learning, ICML, Template}

\vskip 0.3in
]

\printAffiliationAndNotice{}  % leave blank if no need for equal contribution notice

% Abstract
\begin{abstract}
Your abstract should briefly summarize the main contributions of your paper.
It should be a concise paragraph (typically 150-250 words) that covers:
the problem you address, your approach, key results, and implications.
Do not include citations or mathematical notation in the abstract.
\end{abstract}

% Main content sections
\section{Introduction}
\label{sec:introduction}

Introduce your problem and motivate your work. What is the challenge you address?
Why is it important? What gap in existing research does your work fill?

\subsection{Contributions}

List your main contributions:
\begin{itemize}
    \item Contribution 1: Description of first contribution
    \item Contribution 2: Description of second contribution
    \item Contribution 3: Description of third contribution
\end{itemize}

\section{Related Work}
\label{sec:related_work}

Discuss relevant prior work. Position your work relative to existing approaches.
Include citations to important papers in the field \citep{author2024paper}.

\section{Preliminaries}
\label{sec:preliminaries}

Define the problem formally. Introduce necessary notation and background.

\subsection{Problem Formulation}

Let $\mathcal{X}$ be the input space and $\mathcal{Y}$ be the output space.
Given a dataset $\mathcal{D} = \{(x_i, y_i)\}_{i=1}^{n}$, our goal is to learn
a function $f_\theta: \mathcal{X} \rightarrow \mathcal{Y}$ parameterized by $\theta$.

\subsection{Background Concepts}

Explain key concepts that readers need to understand your approach.

\section{Method}
\label{sec:method}

Describe your proposed approach in detail.

\subsection{Overview}

Provide a high-level overview of your method.

\subsection{Detailed Algorithm}

Present the technical details. Use equations where appropriate:

\begin{equation}
    \mathcal{L}(\theta) = \frac{1}{n} \sum_{i=1}^{n} \ell(f_\theta(x_i), y_i)
    \label{eq:loss}
\end{equation}

where $\ell(\cdot, \cdot)$ is the loss function.

\subsection{Algorithm}

\begin{algorithm}[tb]
   \caption{Algorithm Name}
   \label{alg:main}
\begin{algorithmic}
   \STATE {\bfseries Input:} Dataset $\mathcal{D}$, hyperparameters $\eta$, $T$
   \STATE {\bfseries Output:} Model parameters $\theta^*$
   \STATE Initialize $\theta_0$
   \FOR{$t = 1$ {\bfseries to} $T$}
      \STATE Sample mini-batch $\mathcal{B} \subset \mathcal{D}$
      \STATE Compute gradient: $g \leftarrow \nabla_\theta \mathcal{L}(\theta; \mathcal{B})$
      \STATE Update parameters: $\theta \leftarrow \theta - \eta \cdot g$
   \ENDFOR
   \RETURN $\theta^* = \theta$
\end{algorithmic}
\end{algorithm}

\section{Experiments}
\label{sec:experiments}

Describe your experimental setup and results.

\subsection{Datasets}

Describe the datasets used in your experiments.

\subsection{Baselines}

List and describe the baseline methods you compare against.

\subsection{Implementation Details}

Provide implementation details for reproducibility.

\subsection{Results}

Present your main results. Use tables and figures to illustrate findings.

\begin{table}[t]
    \caption{Comparison with baseline methods. Best results are in \textbf{bold}.}
    \label{tab:main_results}
    \vskip 0.15in
    \begin{center}
    \begin{small}
    \begin{sc}
    \begin{tabular}{lcc}
        \toprule
        Method & Accuracy & F1 Score \\
        \midrule
        Baseline 1 & 85.2 & 0.84 \\
        Baseline 2 & 87.3 & 0.86 \\
        Ours & \textbf{91.5} & \textbf{0.90} \\
        \bottomrule
    \end{tabular}
    \end{sc}
    \end{small}
    \end{center}
    \vskip -0.1in
\end{table}

\subsection{Ablation Studies}

Analyze the contribution of different components of your method.

\section{Analysis and Discussion}
\label{sec:discussion}

Discuss your findings, limitations, and broader implications.

\section{Conclusion}
\label{sec:conclusion}

Summarize your contributions and suggest directions for future work.

% Acknowledgments
\section*{Acknowledgments}

This work was supported by [funding sources]. We thank [collaborators] for
helpful discussions.

% Bibliography
\bibliographystyle{icml2025}
\bibliography{references}

% Appendix (if needed)
\appendix
\section{Additional Details}
\label{app:details}

Include additional details in the appendix, such as:
\begin{itemize}
    \item Detailed proofs
    \item Additional experimental results
    \item Hyperparameter sensitivity analysis
\end{itemize}

\end{document}
